\documentclass[12pt]{article}
\usepackage[T1]{fontenc}
\usepackage[utf8]{inputenc}
\usepackage{lmodern}
\usepackage{textcomp}
\usepackage{enumitem}
\usepackage{geometry}
\geometry{margin=1in}
\begin{document}
\begin{center}
Answer Key - Sociology - Graduate Level Assessment 1 \
\small (Answers are ordered exactly as in the question paper.)
\end{center}

\section*{Multiple Choice}
\begin{enumerate}[label=\arabic*.]
\item \textbf{Answer:} D
\\ \textit{Options:} A) ethnic minorities practised religion to achieve social acceptance in the culture; B) mainstream faiths were becoming increasingly identified with national identity; C) the moral teachings of the main religions were becoming relatively similar; D) all of the above
\\ \textit{Note:} Herberg's study highlighted the multifaceted nature of religion in American society, emphasizing that it serves various roles, including as a source of identity, community, and personal meaning, which supports the notion that all presented options collectively reflect his findings.
\item \textbf{Answer:} C
\\ \textit{Options:} A) disappointment and disproportion; B) disbelief and disintegration; C) disengagement and disenchantment; D) distribution and distillation
\\ \textit{Note:} Secularization refers to the process by which religious institutions, practices, and beliefs lose their social significance, encapsulated in the concepts of disengagement from religious authority and disenchantment with religious explanations of the world.
\item \textbf{Answer:} D
\\ \textit{Options:} A) the introduction of compulsory education; B) increasingly emotional ties between parents and children; C) new consumer goods for children, such as clothes, toys and books; D) all of the above
\\ \textit{Note:} The correct answer "all of the above" reflects that the social construction of childhood is influenced by various factors, including historical, cultural, and social changes, which collectively shape our understanding and experience of childhood in different contexts.
\item \textbf{Answer:} C
\\ \textit{Options:} A) popular schools that lay outside their catchment area; B) private, public and comprehensive schools; C) grammar, technical, and secondary modern schools; D) polytechnics, colleges and universities
\\ \textit{Note:} The tripartite system utilized the 11+ exam as a means to classify students into three distinct types of secondary schools-grammar, technical, and secondary modern-based on their academic abilities and potential, reflecting a stratified educational structure aimed at optimizing educational outcomes according to perceived aptitude.
\item \textbf{Answer:} B
\\ \textit{Options:} A) repress it, by promoting only the interests of elite groups; B) revive it, by reaffirming a commitment to freedom of speech; C) reproduce it, by emphasizing face-to-face contact with peer groups; D) replace it with a superior form of communication
\\ \textit{Note:} The correct answer highlights that the Internet has revitalized the public sphere by enhancing opportunities for individuals to express their opinions and engage in dialogue, thereby reinforcing the fundamental democratic principle of freedom of speech.
\end{enumerate}

\section*{True / False}
\begin{enumerate}[label=\arabic*.]
\setcounter{enumi}{5}
\item \textbf{Answer:} False
\\ \textit{Options:} A) True; B) False
\\ \textit{Note:} The nation-state is a fundamental unit of analysis in sociology that encompasses various social, political, and economic structures, often influencing individual and group behaviors more significantly than specific offices within society.
\item \textbf{Answer:} False
\\ \textit{Options:} A) True; B) False
\\ \textit{Note:} Judith Butler argued that gender identity is not solely determined by biological characteristics but is instead a complex interplay of social, cultural, and performative factors, challenging the notion of fixed gender categories based on biology.
\item \textbf{Answer:} True
\\ \textit{Options:} A) True; B) False
\\ \textit{Note:} The social democratic model, as defined by Esping-Andersen, indeed emphasizes universalistic benefits and extensive public sector employment as a means to promote equality and social welfare, reflecting its commitment to providing comprehensive support to all citizens regardless of their socio-economic status.
\item \textbf{Answer:} False
\\ \textit{Options:} A) True; B) False
\\ \textit{Note:} Bell (1973) argued that post-industrial society would prioritize information and services over traditional industrial production and manufacturing, leading to a decline rather than a revival of those practices.
\item \textbf{Answer:} False
\\ \textit{Options:} A) True; B) False
\\ \textit{Note:} The statement is false because race is a socially constructed classification influenced by cultural, historical, and political contexts rather than being solely based on inherent biological traits.
\end{enumerate}

\section*{Long Form Response}
\begin{enumerate}[label=\arabic*.]
\setcounter{enumi}{10}
\item \textbf{Answer:} Warner's findings on the 'colour line' in Natchez illustrate a stark division between black and white castes, reflecting entrenched social hierarchies and pervasive beliefs in white superiority prevalent in the American Deep South. This division is not merely a social construct but is reinforced by institutional practices, cultural norms, and economic disparities that marginalize black communities. Warner highlights how these beliefs justify discriminatory practices and maintain the status quo, perpetuating a cycle of inequality. The rigid boundaries of the colour line serve both to uphold a sense of racial superiority among whites and to systematically dehumanize and disenfranchise black individuals, illustrating the profound impact of these attitudes on societal structure and interpersonal relations.
\\ \textit{Note:} correct because it accurately identifies the 'colour line' as a reflection of systemic racial hierarchies and beliefs in white superiority, emphasizing how these constructs are upheld by institutional practices and economic disparities in the context of the American Deep South.
\item \textbf{Answer:} Social stratification refers to the hierarchical arrangement of individuals or groups within a society based on various factors such as wealth, income, education, and social status. This structure creates distinct social classes that can significantly influence individuals' opportunities and access to resources, leading to varying degrees of social mobility. The implications of social stratification are profound, as they often dictate the life chances of individuals, shaping their experiences in education, employment, and social interactions. For instance, individuals from higher social strata typically benefit from better educational opportunities and professional networks, while those from lower strata may face systemic barriers that limit their mobility and perpetuate inequality. Ultimately, social stratification profoundly affects both individual experiences and broader societal dynamics, reinforcing existing power structures and shaping the collective social landscape.
\\ \textit{Note:} The answer correctly identifies social stratification as a hierarchical system that influences individual opportunities and access to resources, thereby affecting social mobility and shaping personal experiences within society.
\end{enumerate}

\vfill
\noindent\textit{End of Answer Key}
\end{document}